\documentclass[10pt,a4paper]{amsart}
\usepackage[margin=19mm]{geometry}
\usepackage{amsmath,amssymb,mathtools}
\usepackage[scr=rsfs,cal=euler]{mathalfa}
\usepackage{xfrac}
\usepackage[unicode,hidelinks,bookmarks=false]{hyperref}
\newcommand{\ie}{i.e.\ }
\newcommand{\cf}{cf.\ }
\newcommand{\resp}{respectively\ }
\newcommand{\Subsection}{\subsection}
\providecommand{\nobreakdash}{\nobreak\hbox{-}}
\setlength{\emergencystretch}{2em}
\multlinegap0pt
\usepackage{algorithm}\usepackage{algpseudocode}
\usepackage{mathtools}
\usepackage{tikz}
\newtheorem{theorem}{Theorem}
\title{Asynchronous Verifiable Information Dispersal\texorpdfstring{\\}{}with Low Space and Communication Complexity}
\author{Thomas Locher, Yvonne-Anne Pignolet}
\date{Source: arXiv:2608.24636v2 (CC BY 4.0)}
\begin{document}
\maketitle

\noindent\emph{Source and licence.} Asynchronous Verifiable Information Dispersal\texorpdfstring{\\}{}with Low Space and Communication Complexity. Version arXiv:2608.24636v2 (CC BY 4.0), distributed by arXiv under the Creative Commons Attribution 4.0 International licence, http://creativecommons.org/licenses/by/4.0/, as recorded in the arXiv licence metadata for this exact version and shown on the abstract page https://arxiv.org/abs/2608.24636v2. Full licence notice: \texttt{licenses/arxiv-2608.24636-CC-BY-4.0.txt}.\par\smallskip

\noindent\textit{Editorial scope.} Counted unit: the article's theorem theorem (source0.tex lines 788--798). The article's complete original proof of it is delivered with it (source0.tex lines 799--821, one \texttt{proof} environment of 23 lines). No mathematics is written, repaired or re-derived: the statement and the proof are the article's own lines, in the article's own notation, and no retained line is substituted or re-flowed.  Retained uncounted as the source-local context the proof needs: lines 777--780.  Nothing was omitted from any retained span, so the package carries no omission record.  No retained line contains a citation, so the wrapper carries no bibliography and no bibliography entry.  No retained line contains a figure environment, an image-inclusion command or drawing code, so no image is delivered and the package declares no asset dependency.  Editorial formatting only, in addition: an AMS \texttt{amsart} preamble that restates the article's own declarations and macros used by the retained lines, namely the environments \texttt{theorem} on separate counters, together with the standard packages those lines need.  The retained environments print in this document as Theorem 1 in order of appearance, whereas the article prints the same items with the numbering of its own full document.  The complete native source of the article as distributed in the arXiv e-print for this exact version is bundled unchanged as \texttt{sources/208501/source0.tex}.
\begin{theorem}[Space complexity]
\label{theorem:space_complexity}
The space complexity of Protocol $\mathcal{P}$ for storing $m$ is $3|m| + O(n^2|\pi|)$, where $|\pi|$ denotes the size of a commitment proof $\pi$.
\end{theorem}

\begin{theorem}[Communication complexity]
The communication complexity of Protocol $\mathcal{P}$ for $m$ is
%\begin{itemize}
 %\item
 $6|m| + O(n^2|\pi|)$ for dispersal,
 %\item
 $\frac{3}{2}|m| + O(n^2|\pi|)$ for retrieval, and
 %\item
 $\frac{9}{2}\frac{|m|}{n} + O(n|\pi|)$ for recovery.
%\end{itemize}
\end{theorem}

\begin{proof}
\emph{Dispersal}. The client first sends the row and column data to all nodes directly in \emph{disperse} messages, which together require $3|m|+O(n^2|\pi|)$ bits to be sent due to Theorem~\ref{theorem:space_complexity}. The reliable broadcast of $S$ of size $|S| \in O(n|\kappa|)$, where $\kappa$ denotes the size of a signature, has a communication complexity of $O(n^2|\kappa|)$ assuming the use of an efficient reliable broadcast protocol. We further assume that $|\kappa| \in O(|\pi|)$.

Due to $t < n/3$, every element $M[i,j]$ has a size of $$\frac{|m|}{(n-t)^2} < \frac{9}{4}\frac{|m|}{n^2}.$$
The client sends a \emph{spread} message containing $t$ elements, including a proof, to each node. When a node receives a 
\emph{spread} message, it sends two elements and the corresponding proofs, plus the commitment, to each of the $t$ nodes that do not occur in $S$. These two rounds together incur a communication complexity of
$$(1 + 2)\cdot n \cdot t \cdot \left(\frac{9}{4}\frac{|m|}{n^2} + O(|\pi|)\right) < \frac{9}{4}|m| + O(n^2|\pi|).$$

Moreover, each node sends a row element of size $\frac{|m|}{(n-t)^2} < \frac{9}{4}\frac{|m|}{n^2}$ to the $t$ nodes that do not occur in $S$ in a \emph{column\_info} message, which adds
$$n \cdot t \cdot \left(\frac{9}{4}\frac{|m|}{n^2} + O(|\pi|)\right) < \frac{3}{4} |m| + O(n^2|\pi|)$$
to the communication complexity.
Lastly, the constant-sized \emph{ack} and \emph{info\_sent} messages require $O(n)$ and $O(n^2)$ bits to be sent, respectively. Thus, the communication complexity of the entire dispersal process is upper bounded by $6|m| + O(n^2|\pi|)$.
\\
\emph{Retrieval}. Each node sends its stored row data to the client, together with $n-t$ proofs and the commitment, therefore the information sent around amounts to 
\begin{align*}
n\left(\frac{|m|}{(n-t)^2}(n-t) + O(n|\pi|)\right) &= \frac{n|m|}{(n-t)} + O(n^2|\pi|) \\
&< \frac{3}{2}|m| + O(n^2|\pi|),
\end{align*}
again using that $t < n/3$.
\\
\emph{Recovery}. Every node sends two elements of $M$ together with the corresponding proofs and the commitment to the replacement node for a total size of
$$n \cdot 2 \cdot \left(\frac{|m|}{(n-t)^2} + O(|\pi|)\right) < \frac{9}{2}\frac{|m|}{n} + O(n|\pi|).$$
\end{proof}

\end{document}
